
\documentclass[11pt,a4paper]{article}
\usepackage[margin=1in]{geometry}
\usepackage{amsmath,amssymb,amsthm,mathtools}
\usepackage{microtype}
\usepackage{parskip}
\newenvironment{solution}{\par\noindent\textit{Solution.}\ }{\par\vspace{0.5em}}

\begin{document}

\section*{Exercise 3.8 (*)}
Let $h:\mathbb{R}\to\mathbb{R}$ be the $1$-periodic function satisfying $h(x)=|x|$ for $x\in[-\tfrac12,\tfrac12]$
(and $h(x+1)=h(x)$). Define the Takagi function
\[
T(x)=\sum_{i=0}^{\infty}\frac1{2^i}\,h(2^i x).
\]

\begin{enumerate}
\renewcommand\labelenumi{(\alph{enumi})}

\item \textit{Sketch $h$.} On $[-\tfrac12,\tfrac12]$ it is the $V$--shape $|x|$ (cusp at $0$), extended periodically
with period $1$. Equivalently,
\[
h(x)=\mathrm{dist}(x,\mathbb{Z})=\min_{k\in\mathbb{Z}}|x-k|,
\]
so $0\le h(x)\le \tfrac12$, attaining $0$ at integers and $\tfrac12$ at half--integers.

\item \textit{Uniform continuity and differentiability.} Since $x\mapsto \mathrm{dist}(x,\mathbb{Z})$ is $1$--Lipschitz,
$h$ is uniformly continuous on $\mathbb{R}$. On each interval $(k,k+\tfrac12)$ and $(k-\tfrac12,k)$ the function is affine with
slope $-1$ or $+1$, hence differentiable there. It fails to be differentiable precisely at points congruent to
$0$ or $\tfrac12$ modulo $1$, i.e.\ at those $x$ with $2x\in\mathbb{Z}$. Thus $h$ is differentiable at every $x$ with $2x\notin\mathbb{Z}$.

\item \textit{Uniform convergence of the series and uniform continuity of $T$.}
Since $0\le h\le \tfrac12$,
\[
\Big|\frac1{2^i}h(2^i x)\Big|\le \frac{1}{2^{i+1}}\qquad(\forall x\in\mathbb{R}),
\]
and $\sum_{i=0}^\infty 2^{-(i+1)}<\infty$. By the Weierstrass $M$--test, the series converges uniformly on $\mathbb{R}$
to a continuous function $T$. Each summand $x\mapsto 2^{-i}h(2^i x)$ is uniformly continuous (indeed $1$--Lipschitz),
hence $T$, being a uniform limit of uniformly continuous functions, is uniformly continuous on $\mathbb{R}$.

\item Let $x\in\mathbb{R}$. There exist unique $u_n,v_n$ such that
\begin{enumerate}
\renewcommand\labelenumii{(\roman{enumii})}
\item $2^n u_n,\,2^n v_n\in\mathbb{Z}$,
\item $v_n-u_n=2^{-n}$,
\item $u_n\le x< v_n$.
\end{enumerate}
\begin{solution}
Take $u_n=\frac{\lfloor 2^n x\rfloor}{2^n}$ and $v_n=u_n+2^{-n}$. Then (i)--(iii) hold. Uniqueness:
if $u'_n\le x< v'_n$ with (i)--(ii), then $2^n u'_n\le 2^n x<2^n v'_n$ and $2^n v'_n-2^n u'_n=1$, so
$2^n u'_n=\lfloor 2^n x\rfloor=2^n u_n$, hence $u'_n=u_n$ and $v'_n=v_n$.
\end{solution}

\item For $i\ge 0$ set $h_i(x)=2^{-i}h(2^i x)$. Show that:
\begin{enumerate}
\renewcommand\labelenumii{(\roman{enumii})}
\item $h_i(u_n)=h_i(v_n)=0$ for $i\ge n$.
\item $\displaystyle \frac{h_i(v_n)-h_i(u_n)}{v_n-u_n}=\pm 1$ for $i<n$.
\end{enumerate}
\begin{solution}
(i) Since $2^n u_n,2^n v_n\in\mathbb{Z}$, for $i\ge n$ also $2^i u_n,2^i v_n\in\mathbb{Z}$, hence $h(2^i u_n)=h(2^i v_n)=0$.

(ii) The map $x\mapsto 2^i x$ sends $[u_n,v_n]$ onto the interval
$\big[\tfrac{K}{2^{\,n-i}},\tfrac{K+1}{2^{\,n-i}}\big]$ where $K= \lfloor 2^n x\rfloor$.
Its length is $2^{i}(v_n-u_n)=2^{-(n-i)}\le \tfrac12$, so it contains no half--integer.
Hence $h$ is affine with slope $\pm1$ on this interval, and therefore
\[
\frac{h_i(v_n)-h_i(u_n)}{v_n-u_n}=\frac{2^{-i}\big(h(2^i v_n)-h(2^i u_n)\big)}{2^{-n}}
=2^{n-i}\cdot 2^{-i}\cdot (\pm 2^{-(n-i)})=\pm 1.
\]
\end{solution}

\item \textit{Conclude that}
\[
\frac{T(v_n)-T(u_n)}{v_n-u_n}=\sum_{i=0}^{n-1}\frac{h_i(v_n)-h_i(u_n)}{v_n-u_n},
\]
\textit{and that the sum on the right does not converge as $n\to\infty$.}
\begin{solution}
Since $h_i(u_n)=h_i(v_n)=0$ for $i\ge n$,
\[
T(v_n)-T(u_n)=\sum_{i=0}^{\infty}\big(h_i(v_n)-h_i(u_n)\big)
=\sum_{i=0}^{n-1}\big(h_i(v_n)-h_i(u_n)\big).
\]
Divide by $v_n-u_n$ to get the displayed identity. By (e)(ii), each term of the sum equals $\pm1$, so the
sequence of partial sums cannot converge (successive differences are $\pm1\not\to 0$). Hence the right--hand side
does not converge as $n\to\infty$. In particular, $T$ is nowhere differentiable.
\end{solution}
\end{enumerate}

\end{document}

